\documentclass[11pt]{amsart}
\usepackage[margin=1in]{geometry}
\usepackage{amsmath,amssymb,amsthm}
\usepackage{graphicx}
\usepackage[hidelinks]{hyperref}

\newtheorem{theorem}{Theorem}
\newtheorem{lemma}{Lemma}
\newtheorem{corollary}{Corollary}
\theoremstyle{remark}
\newtheorem{remark}{Remark}

\newcommand{\elltwo}{\ell_2}
\newcommand{\N}{\mathbb N}
\newcommand{\J}{\mathcal J}
\newcommand{\F}{\mathcal F}
\newcommand{\Om}{\Omega}
\raggedbottom

\title[Reverse Coordinate Inclusion Obstruction]{A Reverse Coordinate Inclusion Obstruction for a Kalton--Peck Incomparability Conjecture}
\author{agent\_lane\_03}
\date{2026-07-18}

\begin{document}
\begin{abstract}
Castillo and Moreno asked whether, for an interpolation pair $(X^*,X)$ with
a common unconditional basis, a continuous coordinate inclusion $X^*\to X$,
$(X^*,X)_{1/2}=\elltwo$, and singular differential $\Om$, the spaces $X$ and
$X^*$ must be incomparable.  The wording is delicate because one coordinate
inclusion is already assumed.  This note proves a precise partial result: under
the standard complex/CKMR-derived setting used in the source paper, singularity
rules out the reverse coordinate inclusion $X\to X^*$.  Equivalently, the two
endpoint norms cannot be mutually comparable on the common coordinate space.
This does not rule out arbitrary non-coordinate Banach-space embeddings.
\end{abstract}
\maketitle

\section{Source Question}

The source paper is Castillo--Moreno, \emph{The Kalton--Peck space is the
complexification of the real Kalton--Peck space}, arXiv:2204.00672, published in
J. Math. Anal. Appl. 515 (2022), Article~126374.

The relevant passage is in Section 3, ``Concluding remarks'', on page 9 of the
arXiv PDF:

\begin{center}
\includegraphics[width=0.92\textwidth]{figures/open_problem_crop.png}
\end{center}

In text form, the conjecture says:
\begin{quote}
Let $(X,X^*)$ be an interpolation pair of Banach spaces with a common
unconditional basis for which there exists a continuous inclusion $X^*\to X$ and
such that $(X^*,X)_{1/2}=\elltwo$ with differential $\Om$. If $\Om$ is singular
then $X$ and $X^*$ are incomparable.
\end{quote}

Here singularity means that the associated exact sequence has strictly singular
quotient map; equivalently, the restriction of $\Om$ to any infinite-dimensional
subspace is never a sum of a bounded map and a linear map.

\section{Proof Intuition}

The conjecture asks for incomparability, but the hypotheses already include one
coordinate inclusion, $X^*\to X$.  Thus the immediate consistent coordinate
reading is: can the reverse coordinate inclusion $X\to X^*$ also hold?

It cannot.  If both coordinate inclusions are bounded, the two endpoint norms are
equivalent.  An interpolation couple with equivalent endpoints is, up to
equivalent couple norms, just the diagonal couple $(E,E)$.  For a diagonal couple
one can select an interpolation representative by multiplying the vector by a
fixed scalar extremal; the derivative of that selector is a fixed scalar multiple
of the vector.  Hence the differential is linear, modulo the usual bounded
selector-change error.  A globally linear-plus-bounded differential is not
singular.

\section{Main Result}

\begin{theorem}
Let $(X^*,X)$ be a compatible interpolation couple with a common unconditional
coordinate basis and suppose the interpolation method is one of the usual
complex methods, or a CKMR pseudolattice method admitting differentiation, as in
the source paper.  Assume
\[
  (X^*,X)_{1/2}=\elltwo
\]
with associated differential $\Om$, and assume that the coordinate inclusion
$X^*\hookrightarrow X$ is bounded.  If $\Om$ is singular, then the reverse
coordinate inclusion
\[
  X\hookrightarrow X^*
\]
is not bounded.
\end{theorem}

\begin{proof}
We prove the contrapositive.  Suppose that the reverse coordinate inclusion is
also bounded.  Then there are constants $a,b>0$ such that for all finitely
supported coordinate vectors $u$,
\[
  \|u\|_X \le a\|u\|_{X^*}
  \qquad\text{and}\qquad
  \|u\|_{X^*}\le b\|u\|_X .
\]
Thus the two endpoint norms are equivalent on the common coordinate core and the
two completed endpoint spaces can be identified as the same Banach space $E$,
with equivalent norms.

Interpolation is stable under equivalence of compatible couples.  Therefore
$(X^*,X)_{1/2}$ is, through the coordinate identity, equivalent to the diagonal
interpolation space $(E,E)_{1/2}$.  The diagonal interpolation space is just
$E$, with equivalent norm.  Since the hypothesis identifies
$(X^*,X)_{1/2}$ with $\elltwo$, it follows that the endpoint norm of $E$ is
equivalent to the $\elltwo$ norm on the interpolation domain.

It remains to inspect the derived map.  For the diagonal couple $(E,E)$ choose
one fixed scalar admissible representative $\lambda=(\lambda_n)$ for the scalar
unit at the parameter $1/2$; in the complex method this is simply the constant
analytic function $1$, while in the CKMR description one may take any finite
scalar sequence satisfying
\[
  \sum_n e^{n/2}\lambda_n=1 .
\]
For $x\in E$ put $Sx=(\lambda_n x)_n$.  Because the same norm is used at both
endpoints, $S:E\to \J((E,E))$ is bounded.  Its derivative is
\[
  \delta'_{1/2}Sx
   = \left(\sum_n n e^{(n-1)/2}\lambda_n\right)x,
\]
or $0$ in the constant-function complex model.  In either description this is a
bounded linear map on $E\simeq\elltwo$.

Transporting this selector through the endpoint norm equivalence gives a bounded
selector for the original couple $(X^*,X)$ whose differential is bounded linear
on $\elltwo$.  The usual selector-independence statement for derived
interpolation spaces says that replacing one bounded selector by another changes
the differential only by a bounded map into the interpolation space.  Hence the
source differential $\Om$ is a sum of a bounded map and a linear map on all of
$\elltwo$.

This is incompatible with singularity: a singular differential has no
infinite-dimensional restriction which is linear plus bounded.  Therefore the
reverse coordinate inclusion cannot have been bounded.
\end{proof}

\begin{corollary}
Under the hypotheses above, singularity prevents mutual coordinate comparability
of $X$ and $X^*$: the assumed inclusion $X^*\hookrightarrow X$ cannot be
accompanied by a bounded coordinate inclusion $X\hookrightarrow X^*$.
\end{corollary}

\section{Limitations and Novelty Check}

This is a partial result, not a full solution of the broadest possible reading
of the conjecture.  It rules out the reverse coordinate inclusion on the common
basis.  It does not rule out non-coordinate isomorphic embeddings of $X$ into
$X^*$, embeddings of $X^*$ into $X$ with nonclosed range, or stronger notions of
Banach-space incomparability.

Bounded novelty checks were performed before promotion.  The lightweight run
indexes had no exact entry for arXiv:2204.00672.  Web searches for the phrases
``2204.00672 singular differential incomparable'', ``The Kalton-Peck space is
the complexification incomparable'', and ``If Omega is singular then X and X*
are incomparable'' found the source arXiv/JMAA paper and nearby Kalton--Peck
literature, but no later paper explicitly resolving the conjecture.  The 2024
paper \emph{Interpolator Symmetries and New Kalton--Peck Spaces} discusses
related Kalton--Peck interpolation spaces and open questions, but not this
reverse-coordinate obstruction.

\section{Human Review Recommendation}

Review as a likely valid partial theorem.  The key point to check is the
selector-independence step, and the functoriality step, for the exact real or
complex interpolation framework intended by the source conjecture.  The claim
should not be upgraded to a full solution unless ``incomparable'' is explicitly
interpreted as failure of mutual coordinate inclusions.

\begin{thebibliography}{9}

\bibitem{CastilloMoreno}
J. M. F. Castillo and Y. Moreno,
\emph{The Kalton--Peck space is the complexification of the real Kalton--Peck
space},
J. Math. Anal. Appl. 515 (2022), Article 126374.
arXiv:2204.00672.
\url{https://arxiv.org/abs/2204.00672}

\bibitem{CKMR}
M. Cwikel, N. Kalton, M. Milman, and R. Rochberg,
\emph{A unified theory of commutator estimates for a class of interpolation
methods},
Adv. Math. 169 (2002), 241--312.

\bibitem{KaltonPeck}
N. J. Kalton and N. T. Peck,
\emph{Twisted sums of sequence spaces and the three-space problem},
Trans. Amer. Math. Soc. 255 (1979), 1--30.

\bibitem{InterpolatorSymmetries}
J. M. F. Castillo, W. H. G. Correa, V. Ferenczi, and M. Gonzalez,
\emph{Interpolator symmetries and new Kalton--Peck spaces},
Results Math. 79 (2024), Article 108.
\url{https://link.springer.com/article/10.1007/s00025-024-02128-0}

\end{thebibliography}

\end{document}
